\chapter[]{Foundational cryptographic algorithms and techniques}
\section{Introduction}
This chapter will cover the basic algorithms used in cybersecurity, and how they can ensure security properties like confidentiality, authentication or integrity.
These are some terms that must be known:
\begin{description}
    \item[Plaintext ($P$):] The original, readable message.
    \item[Ciphertext ($C$):] The coded (enciphered or encrypted) message.
    \item[Cipher or Algorithm:] The specific technique used to transform plaintext into ciphertext.
    \item[Key:] The essential information used by the cipher to encrypt the plaintext and decrypt the ciphertext.
    \item[Encrypt or Encipher:] The process of converting $P$ to $C$.
    \item[Decrypt or Decipher:] The process of recovering $P$ from $C$.
\end{description}

Encryption is a process used to ensure confidentiality of a message. The encryption algorithms can be divided into 2 categories: symmetric and asymmetric.

\section{Symmetric cryptography}
In symmetric cryptography, there is a single and shared key $K$ between the two parties: the sender who encrypts the data and the receiver who decrypts it. This key must be kept secret. To protect the data, the sender inputs the plaintext message into the encryption algorithm along with the secret key $K$.
The encrypted message (ciphertext) is then transmitted, and the receiver applies the decryption algorithm using the \textbf{same key $K$} to recover the original message.

Note that this mechanism primarily ensures \textit{confidentiality}—meaning only those who possess the key can read the message. It does not inherently guarantee \textit{integrity}. If an attacker intercepts and modifies the ciphertext in transit, the receiver may not be able to detect the tampering (the fact that the message might decrypt into garbled nonsense is not a strict or secure definition of integrity).
A major advantage of symmetric encryption is that it is computationally very fast compared to asymmetric encryption.

Since the key must be shared between parties, a primary challenge that arises is how to distribute it securely. Common methods include:
\begin{itemize}
    \item \textbf{Out of Band (OOB):} The parties meet in person or exchange the key through a completely different, secure communication channel.
    \item \textbf{Trusted Third Party:} Also known as a \textbf{Key Distribution Center (KDC)}. This entity securely issues keys to both parties and is widely used in enterprise protocols like Kerberos.
    \item \textbf{Asymmetric Cryptography:} This is the most common modern approach. It uses slower, public-key cryptography specifically to securely transmit the symmetric secret key to the receiver before switching to the faster symmetric cipher for the actual data.
\end{itemize}

Another critical factor is the length of the key. Kerckhoffs's principle states that a cryptographic system should be secure even if everything about the system, except the key, is public knowledge. Therefore, if:
\begin{itemize}
    \item The key is kept secret and managed only by trusted parties.
    \item The key is sufficiently long.
\end{itemize}
Then the encryption algorithm itself can (and should) be made public. Public algorithms benefit from widespread academic scrutiny, allowing vulnerabilities to be identified and patched. If an algorithm is mathematically sound, the only viable attack is a brute-force attack (exhaustively guessing every possible key). For an $n$-bit key, an attacker must perform up to $2^n$ attempts in the worst-case scenario.

Symmetric algorithms are generally divided into two categories:
\begin{itemize}
    \item \textbf{Block algorithms} (DES, 3DES, AES): These divide the data into fixed-length blocks and process them one at a time. If the plaintext size is not an exact multiple of the block length, padding must be added. When used with the correct modes of operation, the secret key can be safely reused.
    \item \textbf{Stream algorithms} (ChaCha20, Salsa20): These are generally faster than block algorithms and operate on continuous streams of data. They work by taking a master key and expanding it into an infinite, pseudo-random stream of bits called a keystream. The algorithm then encrypts the message by combining the plaintext with the keystream using a bitwise XOR operation.
\end{itemize}

\addfigure[Table of most used symmetric block algorithms (Note: block lengths and key lengths can be required for exercises)]{png/symmetric-algos.png}{0.8}

\subsection{Block algorithms}
The Data Encryption Standard (DES) was one of the first widely adopted symmetric block ciphers. It operates by splitting plaintext into 64-bit blocks and encrypting them using a 56-bit key.
The primary vulnerability of DES lies in this short key length; modern computational power can easily exhaust a 56-bit keyspace via brute-force attacks.

To extend the viability of the algorithm, multiple encryption layers were proposed. An initial approach, Double DES, applied the encryption process twice:
$$C = \text{enc}(K_2, \text{enc}(K_1, P))$$
However, Double DES is critically vulnerable to the meet-in-the-middle (MitM) attack, which drastically reduces its theoretical security and renders it impractical.
Consequently, the industry adopted Triple DES (3DES). 3DES applies the DES algorithm three times in succession, utilizing either two or three distinct 56-bit keys.

The main drawback of 3DES is its slow computational speed, which paved the way for modern standards. Today, the Advanced Encryption Standard (AES) is the globally accepted standard. AES is a block cipher with a fixed block size of 128 bits and supports key lengths of 128, 192, or 256 bits.
The length of the key affects both the computational overhead (longer keys require slightly more processing time) and the lifespan of the data's security (longer keys protect data further into the future). The duration data remains secure is heavily influenced by available computational power; future advancements, particularly in quantum computing, will drastically reduce cracking times. Currently, a 128-bit key is considered secure for decades, while 256-bit keys offer top-secret, quantum-resistant longevity.

\subsection{Modes: Application of block algorithms}
When the data to be encrypted exceeds a single block size, several questions arise: How do we handle a final block that is too short? Should blocks be processed independently, or should they be cryptographically linked?
To address these issues, block algorithms utilize "modes of operation," which define exactly how to securely process multi-block data.

\subsubsection{ECB (Electronic Codebook) mode}
The simplest and most naive mode. The message is divided into blocks, and each block is encrypted entirely independently using the same key.
Formula: $C_i = E_K(P_i)$
\addfigure[ECB encryption and decryption schema]{png/ECB-schema.png}{0.8}
\textbf{Pros:}
\begin{itemize}
    \item \textbf{Parallelization:} Blocks can be encrypted and decrypted in parallel, making the algorithm very fast to execute.
    \item \textbf{No error propagation:} A transmission error in a single ciphertext block generates an error only during the decryption of that specific block; subsequent blocks remain unaffected.
\end{itemize}

\textbf{Cons:}
\begin{itemize}
    \item \textbf{Lack of integrity:} The swapping, reordering, or deletion of ciphertext blocks (e.g., swapping $C_i$ and $C_{i+1}$) goes completely undetected by the receiver.
    \item \textbf{Pattern leakage:} Identical plaintext blocks always generate identical ciphertext blocks under the same key. This reveals data patterns and makes the cipher highly vulnerable to known-plaintext attacks.
\end{itemize}

\textit{Example of a Known-Plaintext Attack on ECB:}
\begin{enumerate}
    \item Assume an attacker computes a database containing the ciphertexts of known-plaintext data (e.g., a highly predictable standard file header) encrypted under many possible keys.
    \item Next, assume the attacker intercepts a ciphertext from the wire or retrieves it from a database.
    \item The attacker can discover the secret key by simply comparing the intercepted ciphertext blocks with the precomputed ciphertexts saved in the database from step 1. Once a match is found, the key is revealed.
\end{enumerate}

\subsubsection{CBC (Cipher Block Chaining) mode} \label{mode:CBC}
To fix the pattern leakage inherent in ECB, CBC introduces chaining. Before being encrypted, each plaintext block is XORed with the previously generated ciphertext block.
Formula: $C_i = E_K(P_i \oplus C_{i-1})$\\
\textbf{Initialization Vector (IV):} Because the very first plaintext block ($P_1$) does not have a "previous ciphertext block," a random IV is used as a stand-in for $C_0$.
\addfigure[CBC encryption schema]{png/CBC-1.png}{0.6}
\addfigure[CBC decryption schema]{png/CBC-2.png}{0.6}

\textbf{Pros:}
\begin{itemize}
    \item \textbf{Resistance to swapping attacks:} Unlike ECB, exchanging two ciphertext blocks (e.g., swapping $C_i$ and $C_{i+1}$) does not go unnoticed. This action will corrupt the corresponding decrypted blocks ($P_i$, $P_{i+1}$, and $P_{i+2}$), making the tampering evident to the receiver.
    \item \textbf{Resistance to known-plaintext attacks:} Because of the chaining mechanism (XORing the current plaintext with the previous ciphertext), identical plaintext blocks will no longer generate identical ciphertext blocks. This effectively hides structural patterns in the data.
\end{itemize}

\textbf{Cons:}
\begin{itemize}
    \item \textbf{Error propagation:} A single transmission error in one ciphertext block (e.g., a flipped bit in $C_i$) will generate decryption errors in exactly two blocks: the current block ($P_i$) will be completely scrambled, and the specific corresponding bit in the next block ($P_{i+1}$) will be flipped.
    \item \textbf{Sequential encryption:} Encryption cannot be parallelized. The algorithm must wait for the computation of $C_i$ to finish before it can begin encrypting $P_{i+1}$. Note: decryption \textit{can} be parallelized because all the cyphertext blocks are available at the beginning.
    \item \textbf{Malleability (Bit-flipping attacks):} CBC mode is malleable. An attacker who knows the structure of the plaintext can strategically invert single bits of a specific plaintext block by flipping the corresponding bits in the Initialization Vector (IV) or the preceding ciphertext block.
\end{itemize}

\subsubsection{Padding}
Padding is not a mode of operation, but rather a structural requirement for many block modes (like ECB and CBC) when the plaintext is not an exact multiple of the block size. For instance, if you need to encrypt 40 bytes of data using a cipher with a 64-byte block size, 24 bytes of padding must be appended.
The exact values used for padding bytes are defined by standardized conventions (such as PKCS\#7).
The two main drawbacks of padding are:
\begin{itemize}
    \item \textbf{Inefficiency:} Padding bytes carry no actual information, meaning you consume bandwidth and storage transmitting useless data.
    \item \textbf{Mandatory Expansion:} Even if the original plaintext is a perfect multiple of the block size, an entire extra block of padding must still be added. Without it, the receiver's decryption algorithm would have no reliable way to distinguish actual data bytes from padding bytes at the end of the file.
\end{itemize}
To avoid the overhead of padding, the following specialized modes and techniques can be used.

\subsubsection{CTS (Ciphertext Stealing)}
CTS is not a standalone mode, but a technique used in conjunction with block cipher modes like CBC or ECB. It allows you to encrypt data that isn't a perfect multiple of the block size \textbf{without adding any padding}. The final ciphertext will be exactly the same length as the original plaintext.
When the algorithm reaches the final, incomplete plaintext block ($P_n$), it "steals" enough bits from the previously generated ciphertext block ($C_{n-1}$) to pad $P_n$ out to a full block. These blocks are then encrypted, and the final two ciphertext blocks are mathematically swapped or truncated to maintain the correct overall message length.
\addfigure[ECB encryption with CTS]{png/CTS-ECB.png}{0.7}
\addfigure[CBC encryption with CTS]{png/CTS-CBC.png}{0.7}

\subsubsection{CTR (Counter) mode}
CTR effectively turns a block cipher into a stream cipher. It eliminates the need for padding and allows data to be encrypted bit-by-bit or byte-by-byte. It uses a \textbf{Counter} (a value that increments for each block) combined with a \textbf{Nonce} (a random Number Used Once).
In this mode, the block algorithm does not encrypt the plaintext directly. Instead, it encrypts the Nonce and Counter to generate a pseudorandom keystream block. This keystream is then XORed with the plaintext to produce the ciphertext.
Formula: $C_i = P_i \oplus E_K(\text{Nonce} \parallel \text{Counter}_i)$

Because the Nonce is required to reconstruct the keystream during decryption, it must be generated by the sender and transmitted (usually in plain text) alongside the ciphertext to the receiver.
\addfigure[CTR encryption]{png/CTR.png}{0.7}

\textbf{Pros:}
\begin{itemize}
    \item \textbf{Parallel execution:} Both encryption and decryption can be highly parallelized.
    \item \textbf{No padding required:} Because the keystream is XORed directly with the plaintext, it operates like a stream cipher and does not require block padding.
    \item \textbf{Random access:} You can directly access and decrypt any specific ciphertext block without needing to process the preceding blocks.
    \item \textbf{Resistance to known-plaintext attacks:} Identical plaintext groups encrypt differently (due to the incrementing counter), making cryptanalysis very difficult.
    \item \textbf{Resistance to swapping attacks:} Exchanging two ciphertext blocks (e.g., $C_i$ and $C_{i+1}$) does not go unnoticed; the resulting decrypted blocks ($P_i$ and $P_{i+1}$) will be corrupted.
\end{itemize}

\textbf{Cons:}
\begin{itemize}
    \item \textbf{Malleability:} A transmission error or intentional bit-flip in one group causes an error \textit{only} in that specific group.
    \item \textbf{Desynchronization:} The cancellation (loss or deletion) of an entire group causes desynchronization between the sender's counter and the receiver's counter. As a result, all subsequent groups will be decrypted erroneously.
\end{itemize}

\begin{table}[h]
    \centering
    \begin{tabular}{|>{\raggedright\arraybackslash}p{3.5cm}|>{\raggedright\arraybackslash}p{3.5cm}|>{\raggedright\arraybackslash}p{3.5cm}|>{\raggedright\arraybackslash}p{3.5cm}|}
        \hline
        \textbf{Feature / Property} & \textbf{ECB} & \textbf{CBC} & \textbf{CTR} \\
        \hline
        \textbf{Parallel Encryption} & Yes (Fast) & No (Sequential bottleneck) & Yes (Fast) \\
        \hline
        \textbf{Parallel Decryption} & Yes & Yes & Yes \\
        \hline
        \textbf{Padding Required} & Yes & Yes & No \\
        \hline
        \textbf{Hides Data Patterns} & No (Identical blocks generate identical ciphertexts) & Yes & Yes \\
        \hline
        \textbf{Transmission Error (1 bit in $C_i$)} & Affects only $P_i$ (No error propagation) & Corrupts $P_i$ and $P_{i+1}$ & Affects only that specific group \\
        \hline
        \textbf{Swapping Attack ($C_i$ and $C_{i+1}$)} & Undetected (Highly vulnerable) & Detected (Corrupts $P_i$, $P_{i+1}$, and $P_{i+2}$) & Detected (Corrupts $P_i$ and $P_{i+1}$) \\
        \hline
        \textbf{Malleability \newline (Bit-flipping attacks)} & Vulnerable to block replacement & Yes (By flipping bits in IV or $C_{i-1}$) & Yes (Highly malleable; direct bit-flipping) \\
        \hline
        \textbf{Random Access to Ciphertext} & Yes & No (Requires previous block to decrypt) & Yes \\
        \hline
        \textbf{Block Cancellation \newline (Dropped block)} & Misaligns data & Misaligns data & Fatal (Desynchronizes counter, corrupting all subsequent groups) \\
        \hline
    \end{tabular}
    \vspace{0.2cm}
    \caption{Comparison of Block Cipher Modes of Operation: ECB, CBC, and CTR}
    \label{tab:modes_comparison}
\end{table}

\subsection{Stream algorithms}
Symmetric stream algorithms operate directly on a continuous stream of data. They encrypt each plaintext digit (typically a single bit or byte) one at a time using a corresponding digit from a keystream, bypassing the need to divide data into fixed-size blocks.
The theoretical ideal for this method is the one-time pad, though it is highly impractical for daily use because it requires a completely random key that is at least as long as the message itself.
To solve this practically, modern stream ciphers utilize pseudo-random key generators to dynamically produce the keystream, which must remain perfectly synchronized between sender and receiver.
Stream algorithms are generally \textbf{faster} than block algorithms and inherently do not require padding. However, they carry a strict security rule: a key/nonce combination must never be reused.
Today, the most prominent and secure stream algorithms are Salsa20 and ChaCha20.

\addfigure[Stream algorithm schema]{png/stream.png}{0.8}


\section{Asymmetric cryptographic}

Asymmetric cryptography uses a linked pair of keys: a public key and a private key. These keys are always generated together and function inversely—a message encrypted with one key can only be decrypted by the other, and vice versa.

Beyond keeping data secret, this method can also be used to ensures authenticity (proving the message truly came from the sender) and integrity (proving the message hasn't been altered).
However, asymmetric encryption is significantly slower than symmetric encryption. Because of this, it is rarely used to encrypt long messages directly.
Instead, it is typically used to securely exchange a faster symmetric key, which is then used to encrypt the actual message data.
Two of the most used algorithm to generate a pair of keys are RSA and Diffie-Hellman

\subsection{RSA}
The RSA (Rivest--Shamir--Adleman) is a foundational public-key algorithm used to generate a pair of public keys.
Its security relies on the mathematical difficulty of factoring the product of two extremely large prime numbers.
The RSA key generation process follows these steps:
\begin{enumerate}
    \item Select two large, secret prime numbers, $P$ and $Q$, randomly chosen.
    \item Compute the public modulus $N = P \times Q$. The overall RSA key size (typically 2048 bits in modern applications) refers to the bit-length of this modulus $N$.
    \item Calculate $\Phi = (P-1)(Q-1)$.
    \item Choose Public Exponent $E$ such that $1 < E < \Phi$, and it is relatively prime to $\Phi$ (i.e. $\gcd(E, \Phi) = 1$).
    \item Compute Private Exponent $D$ as the modular multiplicative inverse of $E$ modulo $\Phi$:
    \[
    D \equiv E^{-1} \pmod{\Phi}
    \]
\end{enumerate}

\noindent This establishes the \textbf{public key} as $(N, E)$ and the \textbf{private key} as $(N, D)$.
During RSA operations, the algorithm can only cipher data where the numerical value of the plaintext $p$ is strictly less than the modulus ($p < N$).

\begin{itemize}
    \item \textbf{Encryption:} $c \equiv p^E \pmod{N}$
    \item \textbf{Decryption:} $p' \equiv c^D \pmod{N}$
\end{itemize}

Terminology Note: Formally, there are no distinct RSA ``encryption'' or ``decryption'' algorithms. There is simply a single RSA mathematical operation. The roles of $E$ and $D$ are entirely interchangeable depending on the key being used, because the math holds true in either direction:
\[
(x^D)^E \pmod{N} \equiv (x^E)^D \pmod{N}
\]

\subsection{Diffie-Hellman}
The Diffie-Hellman (DH) algorithm is classified as key generator, but it is mostly used to securely exchange a shared key, over an insecure, public communication channel.
Unlike RSA, which can be used to directly encrypt messages, Diffie-Hellman is exclusively used to securely generate a symmetric key that both parties can then use to encrypt subsequent communications.

Before the exchange begins, the two communicating parties (traditionally called Alice and Bob) must agree on two public, non-secret parameters:
\begin{enumerate}
    \item \textbf{Prime Modulus ($P$):} A very large prime number.
    \item \textbf{Generator ($G$):} A base number, which is a primitive root modulo $P$.
\end{enumerate}
\noindent \textit{Note: Both $P$ and $G$ can be known to anyone, including potential attackers, without compromising the security of the exchange.}

The process of generating the shared secret happens in the following steps:
\begin{enumerate}
    \item \textbf{Alice's Private and Public Keys:}
    Alice chooses a secret, random integer $a$ (her private key). She then computes her public value $A$ and sends it to Bob:
    \[
    A \equiv G^a \pmod{P}
    \]

    \item \textbf{Bob's Private and Public Keys:}
    Bob chooses his own secret, random integer $b$ (his private key). He computes his public value $B$ and sends it to Alice:
    \[
    B \equiv G^b \pmod{P}
    \]
\end{enumerate}

Once Alice and Bob have exchanged their public values ($A$ and $B$), they independently calculate the final shared secret $S$.
\begin{itemize}
    \item \textbf{Alice computes:} $S \equiv B^a \pmod{P}$
    \item \textbf{Bob computes:} $S \equiv A^b \pmod{P}$
\end{itemize}

Both Alice and Bob arrive at the exact same shared secret $S$ because of the associative property of exponents in modular arithmetic. The mathematical operation holds true from both sides:
\[
S \equiv B^a \pmod{P} \equiv (G^b)^a \pmod{P} \equiv G^{ab} \pmod{P}
\]
\[
S \equiv A^b \pmod{P} \equiv (G^a)^b \pmod{P} \equiv G^{ab} \pmod{P}
\]

\addfigure[DH schema]{png/DH-secure.png}{0.8}

Security Note: An attacker sniffing traffic on the channel only sees $P$, $G$, $A$, and $B$. To figure out the shared secret $S$, they would need to calculate $a$ or $b$.
For sufficiently large values of $P$, calculating the discrete logarithm to find $a$ or $b$ from the public values is computationally infeasible.
However, this basic version of Diffie-Hellman is vulnerable to Man-In-the-middle attacks, because it lacks built-in identity authentication.
When Alice and Bob attempt to exchange their public values over the network, an active attacker intercepts the communication.
Instead of Alice's true public value $A \equiv G^a \pmod{P}$ reaching Bob, the attacker captures it and substitutes it with their own generated public value, $M$.
The attacker does the exact same thing to Bob, intercepting his public value $B \equiv G^b \pmod{P}$ and sending $M$ to Alice instead.
Consequently, Alice unknowingly establishes a shared secret key with the attacker, $S_{Alice} \equiv M^a \pmod{P}$, while Bob establishes a separate shared secret key with the attacker, $S_{Bob} \equiv M^b \pmod{P}$.
Because the attacker holds the corresponding private key for $M$, they can easily compute both $S_{Alice}$ and $S_{Bob}$.
This allows the attacker to silently intercept a message from Alice, decrypt it using $S_{Alice}$, read or modify the plaintext, re-encrypt it using $S_{Bob}$, and forward it to Bob.
Both parties remain completely unaware of the breach, falsely believing they are communicating securely directly with each other.

\addfigure[DH schema]{png/DH-MITM.png}{0.8}

\subsection{Elliptic curve cryptography}
Elliptic curve cryptography is based on the same algorithm that we already saw but instead of using modular arithmetic, the operations are executed on
the surface of a 2D (elliptic) curve.
%<insert advantages and disadvantages>

\section{Message integrity}
So far, we have explored how encryption ensures confidentiality, restricting message access to only the sender and the intended receiver. However, another significant threat arises if an attacker intercepts and modify the message in transit. While decryption might fail or produce unreadable data, relying on the decryption process itself to detect tampering is inefficient. In this chapter, we will examine how a small piece of auxiliary information can be used to verify message integrity \textit{before} decryption. The fundamental approach to achieving this is through the use of a digest, that is a fixed-length summary of a message of arbitrary length (even gigabytes), generated using a cryptographic hash function.

\subsection{Cryptographic Hash Functions}
A cryptographic hash function is an algorithm that accepts an input message of arbitrary length and maps it to a fixed-length output, commonly referred to as a digest. An ideal cryptographic hash function must exhibit the following core properties:
\begin{itemize}
    \item \textbf{Computational Efficiency:} It must be fast to compute the digest for any given message, regardless of its length.
    \item \textbf{Pre-image Resistance (One-Way):} It must be computationally infeasible to invert the function (i.e., to derive the original message from its digest).
    \item \textbf{Collision Resistance:} It must be computationally infeasible to find two distinct messages that result in an identical hash output.
\end{itemize}
Furthermore, in a robust hash function altering even a single bit of the input produces an entirely unpredictable output. Consequently, a hash function serves as a compact, reliable summary of a message, enabling the receiver to verify data integrity efficiently.
Let's see an application of Hash Functions for Integrity Verification.
Consider an object $A$ (e.g., a file, document, or email) whose integrity must be preserved. Initially, the system computes the message digest, $md = h(A)$, and stores this value in a \textbf{secure} location.

Once the digest is secured, the original data $A$ can be stored or transmitted across an insecure environment. To verify the object's integrity at a later time, the system recomputes the hash of the current object and compares it against the securely stored digest $md$. If the two values diverge, it indicates an unauthorized modification, and so a violation of data integrity. Crucially, the digest itself must remain strictly protected; if an adversary gains write access to the digest storage, they could trivially alter the original data, compute a new valid digest, and replace the original, so that the modification goes unnoticed.

\paragraph{MAC, MIC, and MID}
Appending a cryptographic code to a message to guarantee specific security properties is a fundamental practice in security. This methodology relies on three distinctly defined terms:
\begin{itemize}
    \item Message Integrity Code (MIC): A mechanism used exclusively to verify that a message has not been altered in transit, thereby ensuring data integrity.
    \item \textbf{Message Authentication Code (MAC):} A keyed cryptographic checksum used to verify both the integrity of the message and the authenticity of its sender. Because it inherently relies on a shared secret key, it prevents unauthorized parties from generating a valid code.
    \item Message Identifier (MID): A unique value (such as a sequence number, timestamp, or nonce) attached to a message to ensure its uniqueness within a session, effectively protecting the system against replay attacks.
\end{itemize}

\subsection{Keyed Digest}
A keyed digest is a mechanism used to generate a Message Authentication Code (MAC), ensuring both data integrity and origin authentication. It relies on a secret key shared between the sender and the receiver, alongside a mutually agreed cryptographic hash function.

First, the sender computes the digest $md$ over both the message $M$ and the shared secret key $k$. The sender then transmits the message $M$ along with the resulting digest $md$. Then, the receiver computes a new digest $md'$ using the received message $M$ and the same shared key $k$. If the computed digest matches the received digest ($md = md'$), the receiver is assured that the message originated from a party possessing the secret key and that it was not altered in transit.

\addfigure[keyed digest schema]{png/keyed-digest.png}{0.8}

There are two critical limitations to note regarding this approach:
\begin{enumerate}
    \item \textbf{Lack of Non-Repudiation:} Because both the sender and the receiver possess the shared secret key, either party has the necessary information to compute a valid keyed digest. Consequently, the receiver cannot prove to a third party that the sender (and not the receiver themselves) created the message.
    \item \textbf{Data Origin Authentication vs.\ Peer Authentication:} This method authenticates the origin of the specific data packet, but it does not intrinsically authenticate the identity of the peer in a broader session context.
\end{enumerate}

\paragraph{Vulnerabilities and HMAC}
A naive implementation of a keyed digest---such as simply concatenating the key and the message ($h(k \parallel M)$ or $h(M \parallel k)$)---is highly vulnerable to tampering. Many hash functions process data sequentially in blocks. This structure makes them susceptible to \textit{length extension attacks}. If an attacker intercepts the message and the hash, they can append malicious data to the end of the original message and compute a new, valid digest without needing to know the secret key.
To counter these vulnerabilities, the Hash-based Message Authentication Code (\textbf{HMAC}) standard was developed. HMAC securely incorporates the secret key by hashing the data twice using a specific nested construction, effectively neutralizing length extension attacks and providing a robust MAC.

\paragraph{CBC-MAC}
CBC-MAC is another cryptographic mechanism used to generate a MAC. As we've already seen in the \hyperref[mode:CBC]{CBC mode section}, it works by encrypting a message in blocks, where each block's depends on the previous block's ciphertext. This time, the final ciphertext block is then used as the message authentication code.

\addfigure[CBC-MAC schema]{png/CBC-MAC.png}{0.8}

\subsection{Key derivation function}
One the of usage of hash function is to generate secure keys. As the name suggest, a key derivation function (KDF) is a function used to generate a random and secure key that can be used in both symmetric and asymmetric cryptography. The most popular one is password-based key derivation function 2 (PBKDF2), that needs:

\begin{itemize}
    \item password: chosen by the user
    \item salt: that is a random sequence of bytes
    \item hash function
    \item number of iterations: that the hash function is computed
    \item desired key length
\end{itemize}

Using this technique a password chosen by the user can be transformed into a key to be used for cryptography but, the password must be secure.


\subsection{Combining Secrecy and Integrity}
In secure communications, it is often necessary to guarantee both confidentiality (secrecy) and data origin authentication (integrity). Assuming these are handled as distinct operations---using symmetric encryption with a key $K_1$ for secrecy, and a Message Authentication Code (MAC) with a separate key $K_2$ for integrity---there are three primary paradigms for combining them.

It is a fundamental rule in cryptography that improperly combining two secure algorithms can yield an insecure result. The three standard approaches and their security implications are as follows:

\begin{itemize}
    \item \textbf{Authenticate-and-Encrypt (A\&E):} In this approach, the sender computes the ciphertext and the MAC of the plaintext independently, and then concatenates them: $enc(K_1, p) \parallel mac(K_2, p)$.
    \begin{itemize}
        \item \textit{Security Implications:} The receiver is forced to decrypt the message before verifying its integrity. Additionally, computing the MAC directly on the plaintext may inadvertently leak information about the message content. Generally, A\&E is considered insecure unless the operations are tightly coupled in a single step.
    \end{itemize}

    \item \textbf{Authenticate-then-Encrypt (AtE):} Here, the sender first computes the MAC of the plaintext, appends it to the plaintext, and then encrypts the entire package: $enc(K_1, p \parallel mac(K_2, p))$.
    \begin{itemize}
        \item \textit{Security Implications:} Utilized by protocols like SSL and TLS, this method also requires the receiver to perform decryption before the integrity check can occur. Cryptographically, AtE is only proven secure when paired with specific encryption modes, such as Cipher Block Chaining (CBC) or stream ciphers.
    \end{itemize}

    \item \textbf{Encrypt-then-Authenticate (EtA):} In this paradigm, the sender encrypts the plaintext first, and then computes the MAC over the resulting ciphertext: $enc(K_1, p) \parallel mac(K_2, enc(K_1, p))$.
    \begin{itemize}
        \item \textit{Security Implications:} This is widely regarded as the most robust and secure mode, and it is the standard utilized by IPsec. Its primary advantage is that the receiver can evaluate the MAC \textit{before} attempting decryption. If the integrity check fails, the corrupted ciphertext is immediately discarded. However, implementers must still avoid errors; for instance, the Initialization Vector (IV) and algorithm identifiers must always be included in the MAC computation.
    \end{itemize}
\end{itemize}

\paragraph{The Shift to Authenticated Encryption (AE)}
Given the historical pitfalls and complexities of safely combining separate encryption and MAC schemes, modern cryptographic efforts have heavily shifted toward joint \textbf{Authenticated Encryption (AE)} algorithms. These algorithms are specifically designed to provide both confidentiality and authenticity simultaneously, minimizing the risk of implementation errors.

\section{Digital signatures and certificates}
\subsection{Digital signature}
\textbf{Digital signature} is a method similar to the one we saw before to verify integrity with a digest and a symmetric key, but it is performed with a pair of asymmetric keys. The steps are the following:
\begin{itemize}
    \item The sender hashes the data he wants to send, to produce a message digest $md$.
    \item The $md$ undergoes asymmetric encryption using the private key, and the result is a \textbf{digital signature}. This ensures authentication (the key is private, only the sender has it) and integrity (digest).
    \item The receiver receives the data and re-computes the hash $md_R$, while decrypting the received digital signature (asymmetric decryption) to obtain $md_R$.
    \item The receiver compares $md_R$ and $md_S$. If they're the same, then the data has not been modified and it comes from the right sender, because only the sender had that private key.
\end{itemize}
This method achieves integrity but also \textbf{non-repudiation}, because only the sender has that private key, so he can't deny of having created that digital signature.

\addfigure[Digital signature schema]{png/digital-signature.png}{0.8}

\subsection{Public certificate and infrastructure}
In asymmetric cryptography, a fundamental vulnerability is the Man-in-the-Middle (MitM) attack. If an entity, Alice, wishes to communicate securely with Bob, she must obtain Bob's public key. However, if an adversary Eve, intercepts the key exchange and substitutes her own public key for Bob's, she can decrypt, read, and re-encrypt the traffic without detection. Public Key Certificates were introduced to solve this authentication problem by cryptographically \textbf{binding} an entity's identity to their public key, ensuring that the key genuinely belongs to the claimed owner.
A public key certificate functions much like a digital passport. It is a data structure containing an entity's identity details and their public key, which is then digitally signed by a \textbf{trusted third party}, and this is why we need Certification authority.

A Certification Authority (CA) is the trusted third party responsible for verifying identities and issuing public key certificates. Before issuing a certificate, the CA verifies the identity of the applicant. CAs operate in a hierarchy, where a ``Root CA'' signs the certificates of ``Intermediate CAs,'' which in turn sign end-entity certificates. This creates a verifiable \textit{Chain of Trust}.

\paragraph{The X.509v3 Standard}
X.509 is the internationally recognized standard defining the format of public key certificates. Version 3 (X.509v3) is the current and most widely used iteration. It defines a strict ASN.1 data structure containing several mandatory fields:
\begin{itemize}
    \item \textbf{Version:} Indicates the X.509 version (v3).
    \item \textbf{Signature Algorithm Identifier:} The cryptographic algorithm used to sign the certificate (e.g., SHA-256 with RSA).
    \item \textbf{Issuer Name:} The distinguished name of the CA.
    \item \textbf{Validity Period}
    \item \textbf{Subject Name}
    \item \textbf{Subject Email}
\end{itemize}


\paragraph{Revocation Mechanisms}
A certificate is issued with a specific validity period, but it may need to be invalidated before it naturally expires. This occurs if the private key is compromised, if the entity changes its affiliation (e.g., an employee leaves a company), or if the CA itself is compromised. PKI utilizes specific mechanisms to check the revocation status of a certificate. There are two revocation mechanisms:
\begin{itemize}
    \item \textbf{Certificate Revocation List (CRL)}: A CRL is a time-stamped, CA-signed list containing the serial numbers of certificates that have been revoked. It is simple to implement and distribute.
    \item \textbf{Online Certificate Status Protocol (OCSP)}: Instead of downloading an entire list, a client sends an HTTP request to an OCSP responder containing the serial number of the certificate in question. The responder replies with a signed status: \textit{good}, \textit{revoked}, or \textit{unknown}.

\end{itemize}

\paragraph{Example}
The typical scenario for certificates is using OCSP in web to authenticate web servers. Bob, a user, connects to Alice's server in the web. In order for Bob to actually trust Alice, she needs to send her X.509 certificate. But in order for Bob to trust this certificate (that maybe he has never seen before), he wants to know \textbf{who issued it}. This is easy: the entity who issued the certificate is included in the certificate itself. This entity then uses its own certificate to authenticate, which points to another entity who issued it, and so on.

This creates a \textbf{chain} of certificates, that ends when we reach a \textbf{root CA}, a highly trusted public entity that sits at the top of a public key infrastructure (PKI) hierarchy and issues digital certificates. These are the foundation of internet trust, and their certificates are pre-installed in our browsers and operating systems and must be simply trusted.

So, when a certificate is checked, we go up a chain of certificates to verify that the other part is not lying, until we get to a certificate that is self-generated (root CA).
